\section{Practical Constructions}

We have seen private-key schemes constructed from pseudorandom generators,
functions, and permutations. They can be constructed assuming one-way functions exist.
This is a core principle in theoretical modern cryptography.
A one-way function $f : \{0, 1\}^* \rightarrow \{0, 1\}^*$ 
is easy to compute, but hard to invert.
The existence of one-way functions is still an open conjecture.
Assuming one-way functions exist, we can create pseudorandom generators,
functions, and permutations from them.
This is of little practical significance currently, since these constructions are far
less efficient than heuristic methods.
In practice, we use heuristic constructions instead.
Heuristic means they cannot be proven secure based on a weaker assumption.
These constructions have gone through years of public scrutiny and
cryptoanalysis.

\subsection{Stream Ciphers}

Remember our definition of a stream cipher as a 
pair of deterministic algorithms 
(\texttt{Init}, \texttt{Next}):
\texttt{Init} takes input seed $s$; outputs initial state \texttt{st}.
\texttt{Next} takes current state \texttt{st} and
outputs a bit $y$ and updated state \texttt{st'}.
A stream cipher's output behaves like a PRG (indistinguishable from a uniform
sequence).
We can construct a heuristic stream cipher from Linear-Feedback Shift
Registers (LFSR).
However, these have a finite number of states and the 
output periodically repeats. This means that LFSRs are not 
by themselves secure stream ciphers. 
To make them secure, we need to add nonlinearity through 
any of a number of means. 

\subsection{Block Ciphers}

Shannon introduced a paradigm for constructing random-looking permutations.
Construct permutation $F$ with large block length from smaller random (or
random-looking) permutations $\{f_i\}$.
The key for $F$ will describe 16 permutations $f_1, ... , f1_6$ each with block length of
8 bits (1 byte).
Given input $x \in \{0, 1\}^{128}$,
parse as 16 bytes $x_1 ... x_{16}$, then 
set $F_k(x) = f_1(x_1)|| . . . ||f_{16}(x_{16})$.

We say $\{f_i\}$ introduces confusion into $F$.
However, $F$ as described is not pseudorandom. 
If $x$ and $x'$ only differ in first bit, $F_k(x)$ and 
$F_k(x')$ will only differ in first byte.
We add a diffusion step to permute (mix) 
the bits of the output.
The goal is to spread a local change through the entire block.
The confusion and diffusion steps are together called a round.
Rounds are repeated multiple times (also suggested by Shannon).

Substitution-permutation networks (SPN) implement the confusion-diffusion
paradigm.
Permutations $\{f_i\}$ are fixed public substitution 
functions, called an S-box: $f(x) = S(k\oplus x)$.
An important property of an SPN is that a small change in input must affect
every bit of the output.
In SPN, this property translates to a single bit change 
to input of an S-box changes at least two bits in the 
output of the S-box, and mixing permutations ensures 
that bits output by any S-box affect the input to
multiple S-boxes in the next round. 
S-boxes must be designed carefully.

Another approach to construct block ciphers: Feistel Networks.
Each round has a round function $F$ and round key $k$.
Input to round is divided into two halves $L$ and $R$.
Output of the round is two values: $(R, L \oplus F_k(R))$.
To invert, pass output as input to Feistel (with
keys in reverse order).
This works regardless of $F$s, they don't even need
to be invertible.
Three and four round Feistel can be proven to be secure (assuming F is a
pseudorandom function).

\subsection{Hash Functions}

Hash functions map larger input sequence to shorter “digest.”
The main security requirement is collision resistance. 
We want hash functions that can handle very long inputs.
It is much easier to construct fixed-length hash functions (compression functions).
We can extend a fixed-length compression function $h$ into an arbitrary length
input hash function $H$
using the Merkle-Damgard Transform.
We can prove that $H$ is collision resistant if $h$ is collision resistant.
The Merkle-Damgard transform is used in MD5, SHA-1, and SHA-2.
We use the Davies-Meyer construction to build compression functions from block 
ciphers.
Encrypt the prev compression output from the chain using the message block as
key, and XOR output.
Under certain assumptions on the block cipher, the Davies-Meyer construction
can be shown to be collision resistant, as long as $l$ is sufficiently large.

The newest family, SHA-3, moves away from this construction.
SHA-3 uses the sponge construction: data is “absorbed” into the sponge, and
result is “squeezed” out.